\phantomsection\addcontentsline{toc}{section}{Références}
%
\begin{thebibliography}{ABC}
    %
    \bibitem[REF]{reference} auteur. \emph{titre}. édition, année.
    %
    \bibitem[LPP]{lpp} Rolland. \emph{LaTeX par la pratique}. O'Reilly, 1999.
%
\end{thebibliography}




%~ On peut mettre des mots en \emph{italique},
%~ en \textsc{petites Majuscules} ou
%~ en \texttt{largeur fixe (machine à écrire)}.

%~ Voici un deuxième paragraphe avec une formule mathématique simple : $e = mc^2$.

%~ Un troisième avec des \og guillemet français \fg{}.


%~ %
%~ \subsubsection{Écrire en anglais}

%~ \foreignlanguage{english}{Do you speak French? Does anybody here speak french?}


%~ %
%~ \subsection{Listes}

%~ \begin{itemize}
%~ \item Liste classique ;
%~ \item un élément ;
%~ \item et un autre élément.
%~ \end{itemize}
%~ \vspace{\parskip} % espace entre paragraphes

%~ \begin{enumerate}
%~ \item Une liste numéroté
%~ \item deux
%~ \item trois
%~ \end{enumerate}
%~ \vspace{\parskip}

%~ \begin{description}
%~ \item[Description] C'est bien pour des définitions.
%~ \item[Deux] Ou pour faire un liste spéciale.
%~ \end{description}
%~ \vspace{\parskip}


%~ %
%~ \subsection{Références}

%~ Voici une référence à l'image de la figure \ref{latex} page \pageref{latex} et une autre vers la partie \ref{p2} page \pageref{p2}.

%~ On peut citer un livre\,\up{\cite{lpp}} et on précise les détails à la fin du rapport dans la partie références.


%~ %
%~ \subsection{Note de bas de page}

%~ Voici une note\,\footnote{Texte de bas de page} de bas de page.
%~ Une deuxième\,\footnotemark{} déclarée différemment.
%~ La même note\,\footnotemark[\value{footnote}] que précédemment.

%~ \footnotetext{Il a deux références vers cette note}


%~ %
%~ \subsection{Figure}

%~ \begin{figure}[!ht]
    %~ \center
    %~ \includegraphics[]{./images/LaTeX_logo.png}
    %~ \caption{latex | taille original}
    %~ \label{latex}
%~ \end{figure}

%~ \begin{figure}[!ht]
    %~ \center
    %~ \includegraphics[width=0.5\textwidth]{./images/LaTeX_logo.png}
    %~ \caption{latex | 50\% de la largeur de la page}
%~ \end{figure}




%~ ---------------------------------------------------------------------
%~ ---------------------------------------------------------------------
%~ ---------------------------------------------------------------------
